\section{Integer translations and the identity-only symmetry of $Z_1$}
\label{sec:TRZ}
The user's notation is retained: $Z_0$ is absence and zero has no parity; $Z_1$ is presence and no addition on $\tau$ is supplied; $Z_2$ is the even/odd grading of nonzero integers. The user has retracted the earlier additive interpretation of presence. It is not a premise here.

\begin{theorem}[TRZ1: integer addition as translations]
For $m\in\mathbb Z$, put $T_m(n)=n+m$ for all $n\in\mathbb Z$, and set
\[
\operatorname{Trans}(\mathbb Z)=\{T_m:m\in\mathbb Z\}.
\]
This is a subgroup of $\operatorname{Bij}(\mathbb Z)$ under composition, with
\[
T_m\circ T_k=T_{m+k},\qquad T_m^{-1}=T_{-m},\qquad T_0=\operatorname{id}.
\]
The map
\[
\Phi:(\mathbb Z,+)\longrightarrow
(\operatorname{Trans}(\mathbb Z),\circ),\quad m\longmapsto T_m
\]
is an isomorphism, with inverse $T\mapsto T(0)$ on the translation subgroup.
The translation action is faithful, and each nonidentity translation has infinite order.
\end{theorem}
\begin{proof}
For all $m,k,n\in\mathbb Z$,
\[
(T_m\circ T_k)(n)=(n+k)+m=n+(m+k)=T_{m+k}(n).
\]
This uses ordinary integer associativity and commutativity. Taking $k=-m$ in both orders proves the inverse formulas. Thus every $T_m$ is bijective and the family is a subgroup. The composition formula proves that $\Phi$ is a homomorphism; it is surjective by definition. If $T_m=T_k$, evaluation at zero gives $m=T_m(0)=T_k(0)=k$, proving injectivity. Both inverse identities are
\[
\operatorname{ev}_0(\Phi(m))=m,\qquad
\Phi(\operatorname{ev}_0(T_m))=T_m.
\]
A translation acting identically satisfies $m=T_m(0)=0$, so the action is faithful. Induction gives $T_m^r=T_{rm}$ for every positive integer $r$. If $m\ne0$, then $rm\ne0$, hence this is not $T_0$. No parity of zero has been used.
\end{proof}

\begin{theorem}[TRZ2: no faithful factorization through identity-only symmetry]
Let $I_{Z_1}=\{\operatorname{id}\}$ be the group consisting only of the stipulated identity symmetry of $Z_1$. No group homomorphism
\[
\operatorname{Trans}(\mathbb Z)\longrightarrow I_{Z_1}
\]
is injective. Every group homomorphism from $\operatorname{Trans}(\mathbb Z)$ to any group $H$ factoring through $I_{Z_1}$ is trivial, and therefore not faithful.
\end{theorem}
\begin{proof}
Both $T_0$ and $T_1$ have the sole available image, although $T_0(0)=0\ne1=T_1(0)$. A homomorphism from the identity group into $H$ sends its only element to the identity of $H$, so every composite through it is trivial.
\end{proof}
The exact programme conclusion is that the existing identity-only symmetry cannot faithfully represent integer addition through TRZ1. A distinguishable nonidentity translation is additional transformation data. This does not prohibit adding structure externally or imposing a binary operation on a carrier.

\begin{theorem}[TRZ3: parity compatibility with its exact domain]
Put
\[
D=\mathbb Z\setminus\{0\},\quad
D_1=\mathbb Z\setminus\{-1,0\},\quad
D_2=\mathbb Z\setminus\{0,1\}.
\]
Let $\pi:D\to Z_2=\{\mathrm{even},\mathrm{odd}\}$ assign the ordinary parity of each nonzero integer, and let $\sigma$ interchange these two labels. The restriction $T_1:D_1\to D_2$ is a bijection and
\[
\pi\circ(T_1|_{D_1})=\sigma\circ(\pi|_{D_1}).
\]
\end{theorem}
\begin{proof}
Translation by one maps $D_1$ onto $D_2$, with inverse translation by minus one. If $n=2q$, then $n+1=2q+1$ is odd. If $n=2q+1$, then $n+1=2(q+1)$ is even. All inputs and outputs here are nonzero, so no parity at zero is assigned.
\end{proof}
This is a compatibility identity rather than an isomorphism with the parity states: $\pi(1)=\pi(3)$ while $1\ne3$.

\paragraph{TRZ4: the three-state observation.}
On distinct states $0,a,b$, the permutation $R(0)=a$, $R(a)=b$, $R(b)=0$ has $R^3=\operatorname{id}$, whereas $R(0)=a\ne0$ and $R^2(0)=b\ne0$. It therefore already supplies a nonidentity symmetry of order three. It does not derive that symmetry from the identity-only symmetry of $Z_1$, and no parity of the three states is needed.

\paragraph{Scope and provenance.}
TRZ1 begins with integer addition; it does not construct addition from presence or parity. No theorem about all meanings of $\mathbb F_1$, no prismatic comparison, and no claim that every addition requires $Z_2$ is made. The source's lack of addition at $Z_1$ is stipulated, not inferred from the absence of nonidentity symmetries alone.

This is the classical regular action of an additive group, not a new group-theoretic result. For historical human provenance, see Arthur Cayley, \emph{On the theory of groups, as depending on the symbolic equation $\theta^n=1$}, Philosophical Magazine, Series 4, 7(42) (1854), 40--47,
\href{https://doi.org/10.1080/14786445408647421}{doi:10.1080/14786445408647421}.
Publisher metadata was checked; no original author \LaTeX\ was available in that record, and the paper was not read in full. No proof is imported from an unread source: all arguments used here are complete above. The programme framing and notation come from the user's current request; this task wrote the proof and obtained an independent mathematical check.

